\paragraph{\trDirected{}子\trCategory。}\label{sub:directed-subcategories}
暂时地，我们将把自己限制在所有 $\hom$ 空间都是有限维的那一类 \trACategory{}上，
因为这样能简化一些技术细节。设 $\mathbf{Y}=(Y_1,\ldots,Y_m)$ 是这类\trCategory{}%
$\scr{A}$ 中\trObject{}的一个（未必例外的）序列。于是可以定义与之相伴的%
\textbf{\trDirected{}子\trCategory}（directed A∞-subcategory）
$\scr{A}^{\rightarrow}=\scr{A}^{\rightarrow}(\mathbf{Y})$：它是唯一一个以诸
$Y_k$ 为\trObject{}的\trDirected\trACategory，使得当 $i<j$ 时
$\hom_{\scr{A}^{\rightarrow}}(Y_i,Y_j)=\hom_{\scr{A}}(Y_i,Y_j)$，并且其非平庸的%
\trComposition{}映射\eqref{eq:directed-composition-maps} 与 $\scr{A}$ 中相应的
映射一致。这个构造是\trFunctoriality{}的：设 $\scr{F}\colon\scr{A}\to\scr{B}$
是一个 \trAFunctor，并令 $\scr{B}^{\rightarrow}$ 是与像序列 $\scr{F}(\mathbf{Y})$
相伴的\trDirected{}子\trCategory。那么存在唯一一个\trStrictlyUnital{}的%
\trAFunctor{} $\scr{F}^{\rightarrow}\colon\scr{A}^{\rightarrow}\to\scr{B}^{\rightarrow}$，
它在严格上升的\trObject{}链上的限制就等于 $\scr{F}$。若 $\scr{F}$ 是%
\trCohomology\trFullyFaithful{}的，则 $\scr{F}^{\rightarrow}$ 是一个%
\trQuasiisomorphism。

需要提醒的是：若 $\scr{A}$ 是\trStrictlyUnital{}的，则可以把 $\scr{A}^{\rightarrow}$
等同于 $\scr{A}$ 的这样一个子\trCategory：它由诸 $\hom_{\scr{A}}(Y_i,Y_j)$
（$i<j$ 时）以及对每个 $k$ 的子空间 $\KK\cdot e_{Y_k}$ 组成。然而，在一般的%
\trCunital{}情形，$\scr{A}^{\rightarrow}$ 严格来说并不是一个子\trCategory。
为了补救这一点，我们有下面的结果：

\begin{lemma}\label{thm:directed-subcategory-inclusion}
设 $\mathscr{Y}$ 是 $\scr{A}$ 中\trObject{}的一个序列，而 $\scr{A}^{\rightarrow}$
是与之相伴的\trDirected{}子\trCategory。那么存在一个 \trAFunctor{}%
$\mathscr{H}\colon\scr{A}^{\rightarrow}\to\scr{A}$，它在\trObject{}上平凡地
作用，并且对一切 $i<j$ 都诱导\trIsomorphism{}%
$\Hom_{H(\scr{A}^{\rightarrow})}(Y_i,Y_j)\cong\Hom_{H(\scr{A})}(Y_i,Y_j)$。
\end{lemma}

\begin{proof}
利用\cref{thm:strictly-unital-modification} 找到一个\trFormalDiffeomorphism{}%
$\Phi$，使得 $\tilde{\scr{A}}=\Phi_*\scr{A}$ 是\trStrictlyUnital{}的。令
$\tilde{\scr{A}}^{\rightarrow}$ 是 $\tilde{\scr{A}}$ 中与同一序列 $\mathscr{Y}$
相伴的\trDirected{}子\trCategory。由于\trStrictlyUnital{}性，我们有一个包含%
\trFunctor{} $\tilde{\mathscr{H}}\colon\tilde{\scr{A}}^{\rightarrow}
\to\tilde{\scr{A}}$。另一方面，存在一个诱导出的\trFormalDiffeomorphism{}%
$\Phi^{\rightarrow}\colon\scr{A}^{\rightarrow}\to\tilde{\scr{A}}^{\rightarrow}$。
令 $\mathscr{H}=\Phi^{-1}\circ\tilde{\mathscr{H}}\circ\Phi^{\rightarrow}$。
\end{proof}

\begin{lemma}\label{thm:directed-collections-quasi-isomorphic}
设 $\mathbf{Y}$、$\mathbf{Y}'$ 是 $\scr{A}$ 中\trObject{}的两个序列，其性质是
每个 $Y_j$ 都在 $H^0(\scr{A})$ 中与 $Y'_j$ \trIsomorphic。那么与之相伴的%
\trDirected{}子\trCategory{} $\scr{A}^{\rightarrow}(\mathbf{Y})$、
$\scr{A}^{\rightarrow}(\mathbf{Y}')$ 是\trQuasiisomorphic{}的（经由一个保持%
\trObject{}排序的\trStrictlyUnital{} \trAFunctor）。
\end{lemma}

\begin{proof}
令 $\scr{B}\subset\scr{A}$ 是只包含\trObject{} $Y_j$ 的\trFull{} \trAinfty-子%
\trCategory，$\scr{B}'$ 是以诸 $Y'_j$ 为\trObject{}的相应子\trCategory，而
$\scr{C}$ 是同时包含这两类\trObject{}的那一个。按假设，$\scr{B}$ 与 $\scr{B}'$
到 $\scr{C}$ 的包含都是\trQuasiequivalence。于是利用
\cref{thm:quasi-equivalence-homotopy-inverse} 可以构作一个\trQuasiequivalence{}%
$\scr{B}\to\scr{B}'$。从这个构造容易看出，它会把每个 $Y_j$ 映到 $Y'_j$，因而
确实是一个\trQuasiisomorphism。现在过渡到\trDirected{}子\trCategory{}即可。
\end{proof}
\begin{lemma}\label{thm:directed-derived-equivalence}
设 $\mathbf{Y}$ 是一个\trTriangulated{} \trACategory{} $\scr{A}$ 中的一个完整%
\trExceptionalCollection。取 $\scr{A}^{\rightarrow}$ 为与之相伴的\trDirected{}子%
\trCategory，并取 $\mathscr{H}\colon\scr{A}^{\rightarrow}\to\scr{A}$ 为
\cref{thm:directed-subcategory-inclusion} 中的那个 \trAFunctor。那么存在一个%
\trQuasiequivalence{} $\mathscr{G}\colon\mathit{Tw}\scr{A}^{\rightarrow}\to\scr{A}$，
它在 $\scr{A}^{\rightarrow}\subset\mathit{Tw}\scr{A}^{\rightarrow}$ 上的限制在
$H^0(\mathit{fun}(\scr{A}^{\rightarrow},\scr{A}))$ 中与 $\mathscr{H}$
\trIsomorphic。特别地，$D(\scr{A}^{\rightarrow})$ 作为\trTriangulated{}%
\trCategory{}与 $H^0(\scr{A})$ \trEquivalence\footnote{译注：原书此处印作
$H^0(\scr{A}^{\rightarrow})$；疑应为 $H^0(\scr{A})$（按 §3 中对 $D(\cdot)$ 的
定义有 $D(\scr{A}^{\rightarrow})=H^0(\mathit{Tw}\scr{A}^{\rightarrow})$，而本条
前一句恰好构造了\trQuasiequivalence{}%
$\mathscr{G}\colon\mathit{Tw}\scr{A}^{\rightarrow}\to\scr{A}$；且
$H^0(\scr{A}^{\rightarrow})$ 一般并非\trTriangulated{}\trCategory）；此处径改。
——译者注}。
\end{lemma}

\begin{proof}
$\mathbf{Y}$ 是\trExceptionalCollection{}这一事实意味着 $\mathscr{H}$ 是%
\trCohomology\trFullyFaithful{}的。由于 $\scr{A}$ 是\trTriangulated{}的、并且由
$\mathbf{Y}$ 的\trObject{}所生成，它实际上是 $\scr{A}^{\rightarrow}$ 的一个%
\trTriangulated\trEnvelope。对 $\mathit{Tw}\scr{A}^{\rightarrow}$ 而言同样的
结论也成立；现在利用\cref{thm:triangulated-envelope-existence} 即可。
\end{proof}

这个结果的一个应用如下。取一个\trDirected\trACategory{} $\scr{A}^{\rightarrow}$，
以及 $\scr{A}=\mathit{Tw}\scr{A}^{\rightarrow}$ 中由 $\scr{A}^{\rightarrow}$
自身的\trObject{}所构成的完整\trExceptionalCollection。如同\cref{sub:mutation}
中那样，通过施加\trMutation{}来变换这个序列，并令 $\scr{B}^{\rightarrow}$ 是
$\scr{A}$ 中与\trMutation{}后的序列相伴的\trDirected{}子\trCategory。%
\cref{thm:directed-derived-equivalence} 给出一个\trQuasiequivalence{}%
$\mathit{Tw}\scr{B}^{\rightarrow}\to\mathit{Tw}\scr{A}^{\rightarrow}$，从而给出%
\trTriangulated{}\trCategory{}之间的一个\trEquivalence{}%
$D(\scr{B}^{\rightarrow})\cong D(\scr{A}^{\rightarrow})$。正如最初由%
\trKontsevich{}所观察到的那样，这就定义了群 $\ZZ^m\rtimes\mathrm{Br}_m$ 在带
$m$ 个\trObject{}的\trDirected\trACategory{}的\trQuasiisomorphism{}类之集合上
的一个作用，并且具有如下性质：处于同一\trTrajectory{}上的所有\trCategory{}都
具有相互\trEquivalence{}的\trDerivedCategory{}（实际上，中心元素
$(1,\ldots,1)\in\ZZ^m\subset\ZZ^m\rtimes\mathrm{Br}_m$ 平凡地作用，所以该作用
下降到商 $\ZZ^{m-1}\rtimes\mathrm{Br}_m$ 上）。

我们还需要\cref{thm:directed-derived-equivalence}的一个版本，适用于并非%
\trExceptionalCollection{}的序列：

\begin{lemma}\label{thm:spherical-mutation-directed}
设 $\mathbf{Y}$ 是 $\scr{A}$ 中\trObject{}的一个序列，其中每一个都是球性的。在
$\mathit{Tw}\scr{A}$ 中考虑序列 $\mathbf{U}=\mathscr{L}_k\mathbf{Y}$，它由与
\eqref{eq:elementary-mutation} 中相同的公式定义。把与这两个序列分别相伴的%
\trDirected{}子\trCategory{}记作 $\scr{B}^{\rightarrow}$、$\scr{C}^{\rightarrow}$。
那么 $\mathit{Tw}\scr{B}^{\rightarrow}$ 与 $\mathit{Tw}\scr{C}^{\rightarrow}$
是\trQuasiequivalence{}的。
\end{lemma}

\begin{proof}
鉴于\cref{thm:twist-functorial-isomorphism}，我们可以把 \trACategory{} $\scr{A}$
换成任何一个与它\trQuasiequivalence{}的\trCategory，并把 $\mathbf{Y}$、
$\mathbf{U}$ 换成新\trCategory{}中相应的序列。正因为如此，我们可以假设 $\scr{A}$
是\trStrictlyUnital{}的并且 $\mu^1_{\scr{A}}=0$（这纯粹是为了方便）。

我们想把记号稍微区分一下：记 $\tilde{Y}_k$ 为 $\scr{B}^{\rightarrow}$ 中与
$\scr{A}$ 中的 $Y_k$ 相对应的那个\trObject{}。这些\trObject{}构成
$\scr{B}=\mathit{Tw}\scr{B}^{\rightarrow}$ 中的一个完整\trExceptionalCollection{}%
$\widetilde{\mathbf{Y}}$。令
$\widetilde{\mathbf{U}}=\mathscr{L}_k\widetilde{\mathbf{Y}}$\footnote{译注：
原书此处印作 $\widetilde{\mathbf{U}}=\mathscr{L}_k\widetilde{\mathbf{X}}$；疑应为
$\mathscr{L}_k\widetilde{\mathbf{Y}}$（“this collection” 即前句所定义的序列
$\widetilde{\mathbf{Y}}$，而 $\widetilde{\mathbf{X}}$ 在本证明中并无定义）；此处
径改。——译者注}为这个序列的第 $k$ 个初等\trMutation。包含\trFunctor{}%
$\mathscr{H}\colon\scr{B}^{\rightarrow}\to\scr{A}$ 诱导一个 \trAFunctor{}%
$\mathit{Tw}\mathscr{H}\colon\scr{B}\to\mathit{Tw}\scr{A}$，它满足
\begin{equation}\label{eq:tw-h-twist-compatibility}
\mathit{Tw}\mathscr{H}\bigl(T_{\tilde{Y}_i}\tilde{Y}_j\bigr)=T_{Y_i}Y_j
\quad\text{当 } i<j \text{ 时。}
\end{equation}
在左端，我们是在 $\scr{B}$ 的内部进行扭转；而右端则是 $\mathit{Tw}\scr{A}$ 中
相应的构造。两处用的都是\cref{sub:twisted-complexes-representatives}中的显式
公式，于是该等式便化归为
$\hom_{\scr{B}^{\rightarrow}}(\tilde{Y}_i,\tilde{Y}_j)=\hom_{\scr{A}}(Y_i,Y_j)$
这一事实。按定义，\eqref{eq:tw-h-twist-compatibility} 特别地意味着
$\mathit{Tw}\mathscr{H}(\widetilde{\mathbf{U}})=\mathbf{U}$。下一步是要证明映射
\begin{equation}\label{eq:tw-h-first-component}
(\mathit{Tw}\mathscr{H})^1\colon\hom_{\scr{B}}(\widetilde{U}_i,\widetilde{U}_j)
\longrightarrow\hom_{\mathit{Tw}\scr{A}}(U_i,U_j)
\quad\text{当 } i<j \text{ 时}
\end{equation}
都是\trQuasiisomorphism。这只有在 $i$ 或 $j$ 等于 $k$ 时才非平庸。其余的大多数
情形可以这样处理：注意到一端由 $\hom_{\scr{B}^{\rightarrow}}(\tilde{Y}_p,
\tilde{Y}_q)$（$p<q$）搭成，而另一端由相应的 $\hom_{\scr{A}}(Y_p,Y_q)$ 空间
组成；于是按 $\scr{B}^{\rightarrow}$ 的定义，二者重合。唯一剩下的情形是
$(i=k,\ j=k+1)$，此时\eqref{eq:tw-h-first-component} 中的两个群是
\begin{equation*}
\begin{aligned}
\hom_{\scr{B}}\bigl(T_{\tilde{Y}_k}\tilde{Y}_{k+1},\tilde{Y}_k\bigr)
&=\hom_{\scr{A}}(Y_k,Y_{k+1})^{\vee}[-1],\\
\hom_{\mathit{Tw}\scr{A}}(T_{Y_k}Y_{k+1},Y_k)
&=\hom_{\scr{A}}(Y_{k+1},Y_k)\oplus\\
&\qquad\hom_{\scr{A}}(Y_k,Y_{k+1})^{\vee}[-1]\otimes\hom_{\scr{A}}(Y_k,Y_k).
\end{aligned}
\end{equation*}
\begingroup\emergencystretch=8em\hbadness=10000
由于 $\mu^1_{\scr{A}}=0$，第一个群上的\trDifferential{}为零；而在第二个群上，
\trDifferential{}是映射
$(\mu^2_{\scr{A}})'\colon\hom_{\scr{A}}(Y_{k+1},Y_k)\to$
$\hom_{\scr{A}}(Y_k,Y_{k+1})^{\vee}\otimes\hom_{\scr{A}}(Y_k,Y_k)$，
\par\endgroup
它由把\trComposition{}部分地对偶化而得到。球性意味着
\begin{equation*}
(\id\otimes\int)\circ(\mu^2_{\scr{A}})'\colon
\hom_{\scr{A}}(Y_{k+1},Y_k)\longrightarrow
\hom_{\scr{A}}(Y_k,Y_{k+1})^{\vee}[-n_k]
\end{equation*}
是一个\trIsomorphism（$n_k$ 是 $Y_k$ 作为\trSphericalObject{}的维数）。因此
$(\mu^2_{\scr{A}})'$ 是\trInjective{}的，并且由维数计数可知，它的\trImage{}%
是与 $\ker(\id\otimes\int)=\hom_{\scr{A}}(Y_k,Y_{k+1})\otimes\KK e_{Y_k}$ 互补的
一个子空间。但\eqref{eq:tw-h-first-component} 中的映射恰好是
$(\mathit{Tw}\mathscr{H})^1=(0,\id\otimes e_{Y_k})$\footnote{译注：原书此处
印作 $(\mathit{Tw}\mathscr{F})^1$；疑应为 $(\mathit{Tw}\mathscr{H})^1$
（\eqref{eq:tw-h-first-component} 中的映射即 $(\mathit{Tw}\mathscr{H})^1$，
且本证明中并未定义过 $\mathscr{F}$）；此处径改。——译者注}，因而是一个%
\trQuasiisomorphism。

令 $\widetilde{\scr{C}}^{\rightarrow}$ 是 $\scr{B}$ 中与序列
$\widetilde{\mathbf{U}}$ 相伴的\trDirected{}子\trCategory。它带有到 $\scr{B}$
的一个自然\trEmbedding，把它与 $\mathit{Tw}\mathscr{H}$ 组合起来，便得到一个%
\trAFunctor{} $\mathscr{G}\colon\widetilde{\scr{C}}^{\rightarrow}\to
\mathit{Tw}\scr{A}$。这实际上是\trStrictlyUnital{}的；它映到\trDirected{}子%
\trCategory{} $\scr{C}^{\rightarrow}\subset\mathit{Tw}\scr{A}$ 之中；并且由
前面的计算可知，它在 $\widetilde{\scr{C}}^{\rightarrow}$ 与 $\scr{C}^{\rightarrow}$
之间给出一个\trQuasiisomorphism。因此 $\mathit{Tw}\scr{C}^{\rightarrow}$ 与
$\mathit{Tw}\widetilde{\scr{C}}^{\rightarrow}$ \trQuasiequivalence；而由我们
关于\trMutation{}的一般讨论，后者又与 $\mathit{Tw}\scr{B}^{\rightarrow}$
\trQuasiequivalence。
\end{proof}

要把\trDirected{}子\trCategory{}的构造扩充到 $\scr{A}$ 仅为\trCfinite{}的情形，
可以如下进行：利用\trPerturbation{}引理，先过渡到一个与之\trQuasiisomorphic{}%
且实际的 $\hom_{\tilde{\scr{A}}}$ 空间都是有限维的 \trACategory{}%
$\tilde{\scr{A}}$，然后把前面的构造应用到那个\trCategory{}上。所得结果仍记作
$\scr{A}^{\rightarrow}$，它再次带有一个如\cref{thm:directed-subcategory-inclusion}
中的 \trAFunctor{} $\mathscr{H}\colon\scr{A}^{\rightarrow}\to\scr{A}$。缺点是
$\scr{A}^{\rightarrow}$ 不再是\trCanonical{}的，因为它依赖于 $\tilde{\scr{A}}$
的选取；但它至少在相差\trQuasiisomorphism{}的意义下是唯一的，这对大多数目的
而言已经足够。特别地，\cref{thm:directed-derived-equivalence}与%
\cref{thm:spherical-mutation-directed} 将照旧成立。
